\chapter{СИСТЕМИЙН КОДЫН ХАВСРАЛТ}

Энэхүү хавсралтад системийн гол бүрэлдэхүүн хэсгүүдийн кодыг төлөөлөх хэлбэрээр багтаав. Хавсралтад бүх эх кодыг хуулбарлан оруулахын оронд хамгаалалтын үед тайлбарлахад шаардлагатай үндсэн сервис, API, клиент, өгөгдлийн сангийн холболтын хэсгүүдийг сонгон үзүүлсэн.

\begingroup
\lstset{
    style=mystyle,
    basicstyle=\ttfamily\footnotesize,
    numbersep=6pt,
    aboveskip=6pt,
    belowskip=6pt
}

\section{Аудитын сервисийн үндсэн логик}

Дараах код нь хүлээгдэж буй болон илрүүлсэн барааны тоог харьцуулж \texttt{PASS}, \texttt{WARNING}, \texttt{FAIL} статус гаргах гол логикийг харуулна.

\begin{lstlisting}[language=Python, caption=AuditService-ийн үндсэн аудитын функц]
class AuditService:
    def __init__(self, tolerance: float = 0.1,
                 critical_threshold: float = 0.3):
        self.tolerance = tolerance
        self.critical_threshold = critical_threshold

    def run_audit(self, expected: Dict[str, int],
                  detected: Dict[str, int]) -> AuditResult:
        discrepancies = []
        missing_products = []
        extra_products = []
        total_expected = sum(expected.values())
        total_matched = 0

        for item, count in expected.items():
            found = detected.get(item, 0)
            diff = found - count

            if found == 0:
                missing_products.append(item)
            elif diff != 0:
                discrepancies.append(ProductDiscrepancy(
                    product_name=item,
                    expected_count=count,
                    detected_count=found,
                    difference=diff
                ))

            total_matched += min(found, count)

        for item in detected:
            if item not in expected:
                extra_products.append(item)

        match_rate = 0.0 if total_expected == 0 else \
            (total_matched / total_expected) * 100

        if match_rate >= 100:
            status = AuditStatus.PASS
        elif match_rate >= 80:
            status = AuditStatus.WARNING
        else:
            status = AuditStatus.FAIL

        return AuditResult(
            status=status,
            match_rate=round(match_rate, 2),
            discrepancies=discrepancies,
            missing_products=missing_products,
            extra_products=extra_products,
            total_expected=total_expected,
            total_detected=sum(detected.values())
        )
\end{lstlisting}

\section{Бараа илрүүлэх сервисийн төлөөлөх хэрэгжүүлэлт}

Энэ хэсэгт Singleton загвараар загварыг нэг удаа санах ойд ачаалж, дараагийн хүсэлтүүдэд дахин ашиглах байдлыг харуулав. Бүрэн хэрэгжүүлэлтэд алдааны боловсруулалт, тохиргооны уншилт, зураг тэмдэглэгээ зэрэг нэмэлт функцууд багтсан.

\begin{lstlisting}[language=Python, caption=ProductDetectionService-ийн гол бүтэц]
from ultralytics import YOLO

class ProductDetectionService:
    _instance = None

    def __new__(cls):
        if cls._instance is None:
            cls._instance = super().__new__(cls)
            cls._instance._initialize()
        return cls._instance

    def _initialize(self):
        self.model = None
        self.confidence_threshold = 0.25
        self.iou_threshold = 0.45
        self.class_names = {}

    def load_model(self, model_path: str):
        self.model = YOLO(model_path)
        self.class_names = self.model.names

    def detect(self, image: np.ndarray) -> DetectionResult:
        if self.model is None:
            raise ValueError("Model not loaded")

        results = self.model(
            image,
            conf=self.confidence_threshold,
            iou=self.iou_threshold,
            verbose=False
        )[0]

        detections = []
        for box in results.boxes:
            class_id = int(box.cls)
            detections.append(Detection(
                class_id=class_id,
                class_name=self.class_names[class_id],
                confidence=float(box.conf),
                bbox=box.xyxy[0].tolist()
            ))

        return DetectionResult(
            detections=detections,
            total_products=len(detections),
            processing_time_ms=0.0
        )
\end{lstlisting}

\section{FastAPI сервисийн endpoint-ийн төлөөлөх хэсэг}

Энэ endpoint нь мобайл эсвэл вэб клиентаас ирсэн зургийг хүлээн авч, inference хийж, үр дүнг MongoDB-д хадгалж буцаах урсгалыг нэг дороос харуулж байна.

\begin{lstlisting}[language=Python, caption=Detection API сервисийн гол урсгал]
@router.post("/detect")
async def detect_products(file: UploadFile = File(...)):
    if not file.content_type.startswith("image/"):
        raise HTTPException(status_code=400, detail="Invalid file type.")

    contents = await file.read()
    nparr = np.frombuffer(contents, np.uint8)
    image = cv2.imdecode(nparr, cv2.IMREAD_COLOR)

    if image is None:
        raise HTTPException(status_code=400, detail="Could not decode image")

    result = detection_service.detect(image)
    collection = await get_detections_collection()

    doc = {
        "timestamp": datetime.utcnow(),
        "filename": file.filename,
        "detections": [
            {
                "class_name": d.class_name,
                "confidence": d.confidence,
                "bbox": d.bbox
            }
            for d in result.detections
        ],
        "total_products": result.total_products,
        "processing_time_ms": result.processing_time_ms
    }

    insert_result = await collection.insert_one(doc)

    return JSONResponse({
        "success": True,
        "detection_id": str(insert_result.inserted_id),
        "total_products": result.total_products,
        "processing_time_ms": result.processing_time_ms,
        "detections": doc["detections"]
    })
\end{lstlisting}

\section{Flutter клиент талын API сервис}

Мобайл апп талд зургийг \texttt{multipart/form-data} хэлбэрээр илгээж, серверийн JSON хариуг \texttt{DetectionResponse} обьект руу хөрвүүлэн ашигладаг.

\begin{lstlisting}[language=Java, caption=Flutter API сервисийн төлөөлөх функц]
class ApiService {
  static const String baseUrl = 'http://localhost:8000/api';

  static Future<DetectionResponse> detectProducts(
    File imageFile
  ) async {
    var request = http.MultipartRequest(
      'POST',
      Uri.parse('$baseUrl/detection/detect'),
    );

    request.files.add(
      await http.MultipartFile.fromPath('file', imageFile.path),
    );

    var streamedResponse = await request.send();
    var response = await http.Response.fromStream(streamedResponse);

    if (response.statusCode == 200) {
      return DetectionResponse.fromJson(json.decode(response.body));
    }

    throw ApiException('Detection failed: ${response.body}');
  }
}
\end{lstlisting}

\section{MongoDB холболтын сервис}

Өгөгдлийн сангийн бүрэн тохиргоо, collection accessor-уудыг багасган үзүүлэв. Үндсэн санаа нь FastAPI-ийн амьдралын мөчлөг дотор MongoDB-тэй холбогдож, асинхрон driver ашиглан collection-уудыг авахад оршино.

\begin{lstlisting}[language=Python, caption=Motor async driver ашигласан өгөгдлийн сангийн сервис]
from motor.motor_asyncio import AsyncIOMotorClient

class Database:
    client: AsyncIOMotorClient = None

db = Database()

async def connect_db():
    db.client = AsyncIOMotorClient(settings.MONGODB_URI)
    await db.client.admin.command('ping')

async def disconnect_db():
    if db.client:
        db.client.close()

def get_database():
    return db.client[settings.MONGODB_DB_NAME]
\end{lstlisting}

\section{Хавсралтын бүтэц сонгосон үндэслэл}

\begin{itemize}
    \item Хавсралтад бүрэн эх код бус, архитектурын хувьд хамгийн төлөөлөх хэсгүүдийг үлдээв.
    \item Давтагдсан boilerplate, DTO, import, helper функцуудыг хасч, хамгаалалтын үед тайлбарлах шаардлагатай үндсэн логикийг хадгалав.
    \item Ингэснээр хавсралт мэргэжлийн хэв шинжээ алдахгүйгээр илүү нягт, уншигдахуйц болж, тайлангийн нийт хуудасны тоог оновчтой түвшинд барих боломж бүрдэнэ.
\end{itemize}

\endgroup
